\chapter{줄 세우기 비교}

엔도스랭크와 송금 랭크로 지갑 5,521개에 점수를 매겼어요.
이제 궁금한 것이 생겨요.
두 점수는 얼마나 비슷할까요?
비슷하다면 굳이 새 방법을 만들 필요가 없겠지요.

\section{두 줄이 얼마나 닮았나}

체육 시간에 다섯 친구가 줄을 섰어요.
한 번은 키 큰 순서로, 한 번은 나이 많은 순서로요.

\begin{figure}[h]
\centering
\begin{tikzpicture}[x=1.9cm]
  \node[font=\small, anchor=east] at (-0.2,1.6) {키 순서};
  \node[font=\small, anchor=east] at (-0.2,0) {나이 순서};
  \foreach \i/\k/\n in {0/a/지우, 1/b/하준, 2/c/서연, 3/d/도윤, 4/e/민서}
    \node[kid, minimum size=9mm, font=\scriptsize] (t\k) at (\i,1.6) {\n};
  \foreach \i/\k/\n in {0/b/하준, 1/a/지우, 2/c/서연, 3/e/민서, 4/d/도윤}
    \node[kid, minimum size=9mm, font=\scriptsize] (b\k) at (\i,0) {\n};
  \draw[line width=1pt, color=leafline] (tc) -- (bc);
  \foreach \k in {a,b,d,e}
    \draw[line width=1pt, color=roseline] (t\k) -- (b\k);
\end{tikzpicture}
\caption{같은 친구끼리 선으로 이었어요. 선이 엇갈린 곳이 두 군데 있어요.}
\end{figure}

두 줄을 비교하는 방법은 이래요.
다섯 명 가운데 두 명씩 짝을 지으면 모두 열 쌍이 나와요.
쌍마다 이렇게 물어요.
“키 순서에서 앞에 선 친구가 나이 순서에서도 앞에 섰나?”

\begin{itemize}
\item 그렇다면 \textbf{맞는 쌍}이에요. 지우와 서연이는 두 줄 모두 지우가 앞이니 맞는 쌍이에요.
\item 아니라면 \textbf{어긋난 쌍}이에요. 지우와 하준이는 키로는 지우가 앞인데 나이로는 하준이가 앞이에요.
\end{itemize}

열 쌍을 다 따져 보면 어긋난 쌍은 (지우, 하준)과 (도윤, 민서) 둘뿐이고, 나머지 여덟 쌍은 맞아요.
그림에서 선이 엇갈린 곳의 수가 바로 어긋난 쌍의 수예요.

이제 점수를 만들어요.
\[
\text{닮은 정도} = \frac{\text{맞는 쌍} - \text{어긋난 쌍}}{\text{모든 쌍}} = \frac{8-2}{10} = 0.6
\]

이 점수를 \textbf{켄달의 타우}라고 부르고, 그리스 문자 $\tau$로 써요.
1938년 영국의 통계학자 모리스 켄달이 만들었어요.
두 줄이 완전히 같으면 1, 완전히 거꾸로면 $-1$, 아무 관계가 없으면 0 근처가 나와요.
이 책에서 앞으로 나오는 비교는 거의 모두 이 $\tau$로 해요.

\section{엔도스랭크와 송금 랭크는 얼마나 닮았을까?}

지갑 5,521개를 엔도스랭크 점수 순서로 한 줄, 송금 랭크 점수 순서로 한 줄 세웠어요.
두 줄의 $\tau$는 \textbf{0.360}이었어요.

1이 아니에요.
0.6도 아니에요.
두 점수는 조금 닮았지만, 지갑을 꽤 다르게 줄 세운다는 뜻이에요.
같은 지갑들을 보고도 “누가 믿을 만한가”에 대해 다른 대답을 하는 거예요.
이것이 논문의 첫 번째 질문에 대한 답이에요.

\begin{deeper}[0.360을 조금 더 들여다보기]
그나마 닮은 부분도 상당 부분은 허락 지도에 나타나지 않아 엔도스랭크 점수가 0인 지갑 966개에서 나와요.
이 지갑들은 모두 같은 점수(0)를 받아 줄의 뒤쪽에 함께 서는데, 이들은 송금 랭크에서도 점수가 낮은 편이에요.
그래서 논문에서는, 지도 밖의 지갑을 점수 0 대신 “화살표 없이 지도에 들어 있는 점”으로 다루어 다시 계산해 보았어요.
그러자 두 점수의 $\tau$는 $-0.025$, 거의 0이 되었어요.
송금 랭크는 송금을 보낸 적 있는 지갑에서만 다시 출발하기 때문에, 이렇게 다루어도 송금 지도 밖 지갑의 점수는 0 그대로예요.
이 규칙은 결과를 본 뒤에 시도한 것이라 참고로만 보고해요.
\end{deeper}

\section{어느 쪽이 무엇과 닮았나}

두 점수가 다르다는 건 알았어요.
그럼 각각은 무엇을 잘 나타낼까요?
이것을 알아보려고 지갑마다 아주 단순한 사실들을 세었어요.
이런 사실들을 \textbf{검증 지표}라고 불러요.

\begin{itemize}
\item \textcolor{sendcol}{\textbf{송금 검사}}: 송금을 몇 번 받았는지, 얼마나 받았는지
\item \textcolor{allowcol}{\textbf{허락 검사}}: 허락을 몇 번 받았는지, 얼마나 받았는지
\end{itemize}

그리고 점수 줄과 검사 줄의 $\tau$를 쟀어요.
결과는 이랬어요.

\begin{figure}[h]
\centering
\begin{tikzpicture}
\begin{axis}[
  ybar, bar width=12pt, width=0.95\linewidth, height=5.4cm,
  ymin=-0.1, ymax=0.72, ytick={0,0.2,0.4,0.6},
  symbolic x coords={A, T},
  xtick=data, xticklabels={허락 검사, 송금 검사}, enlarge x limits=0.5,
  legend style={font=\scriptsize, at={(0.5,1.02)}, anchor=south, legend columns=2, draw=none},
  nodes near coords, nodes near coords style={font=\scriptsize},
  every node near coord/.append style={/pgf/number format/fixed, /pgf/number format/fixed zerofill, /pgf/number format/precision=3},
  tick label style={font=\small}, ylabel={$\tau$}, axis x line*=bottom,
]
\addplot[fill=allowcol, draw=allowcol] coordinates {(A,0.375) (T,0.330)};
\addplot[fill=sendcol, draw=sendcol] coordinates {(A,-0.024) (T,0.612)};
\legend{엔도스랭크, 송금 랭크}
\end{axis}
\end{tikzpicture}
\caption{허락 검사에서는 엔도스랭크가, 송금 검사에서는 송금 랭크가 더 닮았어요.}
\end{figure}

허락 검사에서는 엔도스랭크가 0.375, 송금 랭크가 $-0.024$였어요.
송금 검사에서는 송금 랭크가 0.612, 엔도스랭크가 0.330이었어요.
각자 자기가 쓴 기록과 닮은 쪽에서 더 강했어요.
달리기 1등과 받아쓰기 1등이 다른 것처럼요.
두 점수는 서로 다른 종류의 믿음을 재고 있다는 뜻이에요.

\section{우연은 아닐까?}

그런데 이런 차이가 우연히 나온 것일 수도 있잖아요?
그래서 과학자들은 \textbf{다시 뽑기}를 해 봐요.

지갑 5,521개가 적힌 쪽지를 상자에 넣고, 한 장 뽑아 적고 다시 넣기를 5,521번 해요.
그러면 어떤 지갑은 두 번 뽑히고 어떤 지갑은 안 뽑힌, 조금 다른 무리가 생겨요.
그 무리로 $\tau$를 다시 계산해요.
이것을 400번 되풀이하면 $\tau$가 400개 나와요.
그 가운데 가장 작은 쪽 2.5\%와 가장 큰 쪽 2.5\%를 빼고 남은 범위를 \textbf{95\% 범위}라고 해요.
“운이 조금 나쁘거나 조금 좋아도 대개 이 안에 들어온다”는 범위예요.

엔도스랭크와 송금 랭크의 $\tau$ 0.360은 95\% 범위가 0.342에서 0.379였어요.
범위가 좁고 1에서 멀리 떨어져 있으니, 두 점수가 다르다는 건 우연이 아니에요.
허락 검사와 송금 검사의 차이도 같은 방법으로 확인했어요.
이 다시 뽑기는 14장에서 다시 아주 중요해져요.

\begin{deeper}[진짜 이름과 숫자]
\begin{itemize}
\item 논문에서는 동점을 고려한 켄달의 $\tau_b$를 썼어요.
\item 다시 뽑기는 부트스트랩이라고 해요. 같은 지갑 표본으로 두 방법을 함께 계산하는 “짝지은 부트스트랩”을 400번, 시드 42로 했어요.
\item 검증 지표는 여섯 가족으로 묶었어요: 송금, 허락, GMX 청산, 손실 위험, 시빌 안정성, GMX 거래 성과. 본문의 주된 비교는 송금과 허락 가족이에요. 시빌 안정성은 송금 랭크 0.278, 엔도스랭크 0.231이었고, GMX 거래 성과와는 둘 다 약하게만 닮았어요(송금 랭크 0.236, 엔도스랭크 0.096).
\item 두 검사의 가장 높은 가능 값이 서로 달라서, “엔도스랭크가 송금 검사보다 허락 검사와 더 닮았다”고 말하지는 않아요. 같은 검사 안에서 두 방법을 비교할 뿐이에요.
\item 엔도스랭크와 송금 랭크 점수는 이 검사들과 봄 채점(13장)의 결과가 나온 뒤에 매겼어요. 그래서 이 장과 13장에서 두 방법을 맞대어 본 비교는 미리 정해 둔 것이 아니고, 범위와 함께 그대로 보고해요.
\end{itemize}
\end{deeper}

\begin{learned}
\begin{itemize}
\item 켄달의 $\tau$는 두 줄에서 맞는 쌍과 어긋난 쌍을 세어 두 순서가 얼마나 닮았는지 재요.
\item 엔도스랭크와 송금 랭크의 $\tau$는 0.360이에요. 두 점수는 지갑을 꽤 다르게 줄 세워요.
\item 허락 검사에서는 엔도스랭크가, 송금 검사에서는 송금 랭크가 더 닮았어요.
\item 다시 뽑기를 400번 해서 차이가 우연이 아닌지 확인했어요.
\end{itemize}
\end{learned}
